\section{Flatten a Linked List}
\label{sec:ll-flatten}

\problemstatement{Each node of a linked list has a \texttt{next} pointer (to
the next sub-list head) and a \texttt{bottom} pointer heading a sorted
vertical sub-list. Flatten the whole structure into a single sorted list
linked entirely by \texttt{bottom} pointers.}

\begin{patternbox}
Every vertical sub-list is already sorted, so flattening is \textbf{repeated
merging}: merge two sorted \texttt{bottom}-lists at a time, folding the
horizontal \texttt{next} chain into one long sorted vertical list. It is the
merge routine of merge sort applied across sub-lists.
\end{patternbox}

\subsection*{Merge sub-lists right to left}

Recurse along \texttt{next} to flatten everything to the right of the head
first, then merge the head's vertical list with that flattened remainder.
The merge compares \texttt{bottom} values and links the smaller each time,
producing one sorted \texttt{bottom}-chain. Processing right-to-left means
each merge combines the current sub-list with an already-fully-flattened
tail.

\begin{ideabox}[Fold the Horizontal Lists into One Vertical Merge]
Because each vertical list is sorted, the whole grid flattens by pairwise
merging along the \texttt{next} direction. The result uses only
\texttt{bottom} pointers, discarding the horizontal structure.
\end{ideabox}

\subsection*{Algorithm}

\begin{enumerate}
  \item If \texttt{head} or \texttt{head->next} is null, return
        \texttt{head}.
  \item Recursively flatten \texttt{head->next}.
  \item Merge \texttt{head}'s vertical list with the flattened remainder by
        \texttt{bottom} value; return the merged head.
\end{enumerate}

\subsection*{Dry run}

\dryrun{Two sub-lists: $3\!\downarrow\!7$ and $1\!\downarrow\!4$ (down =
bottom), linked by \texttt{next}.}

\begin{itemize}
  \item Flatten the second: $1\to4$. Merge with $3\to7$.
  \item Result (bottom-linked) $1\to3\to4\to7$.
\end{itemize}

\subsection*{C++ implementation}

\begin{lstlisting}
#include <bits/stdc++.h>
using namespace std;

struct Node {
    int val;
    Node* next;                                   // horizontal
    Node* bottom;                                 // vertical (sorted)
    Node(int v) : val(v), next(nullptr), bottom(nullptr) {}
};

Node* merge(Node* a, Node* b) {                   // merge two bottom-lists
    Node dummy(0), *t = &dummy;
    while (a && b) {
        if (a->val <= b->val) { t->bottom = a; a = a->bottom; }
        else                  { t->bottom = b; b = b->bottom; }
        t = t->bottom;
    }
    t->bottom = a ? a : b;
    return dummy.bottom;
}

Node* flatten(Node* head) {
    if (!head || !head->next) return head;
    head->next = flatten(head->next);             // flatten the rest first
    head = merge(head, head->next);
    return head;
}

int main() {
    Node* a = new Node(3); a->bottom = new Node(7);
    Node* b = new Node(1); b->bottom = new Node(4);
    a->next = b;
    Node* f = flatten(a);
    for (Node* c = f; c; c = c->bottom) cout << c->val << " ";
    cout << "\n";                                  // 1 3 4 7
    return 0;
}
\end{lstlisting}

\begin{complexitybox}
\textbf{Time Complexity.} $O(N)$ total merge work across all
$N$ nodes (with $O(N\cdot m)$ in the naive per-merge bound, $m$ sub-lists).
\textbf{Space Complexity.} $O(m)$ recursion over the \texttt{next} chain.
\end{complexitybox}

\begin{takeawaybox}
Flattening sorted sub-lists is repeated merging along the horizontal
direction, reusing the merge-sort combine step. Recognising ``already
sorted pieces $\Rightarrow$ merge, don't re-sort'' is the key insight; a
priority queue over sub-list heads is an alternative.
\end{takeawaybox}
